% ===== ch02 : 原文 main.tex 第 365--515 行（2 定理的证明）=====
% 本文件是原文 LaTeX 源码的逐行分片（去掉整段注释掉的旧稿），仅供翻译时对照；最终译文写到 src/chunks/ch02.tex。
\section{Proof of the Theorems\label{sec:proof-of-theorems}}
In this section, with technical details postponed to later sections, we will provide proofs of the claimed results in the Introduction.
\subsection{The preliminary and the blueprint.} Throughout this paper, on the $A$-model we will regard all symplectic manifolds $X$ as \emph{Liouville sectors}, which will be defined in Definition \ref{defn:Liouville-sector-inf}. For Liouville sectors of finite type their wrapped Fukaya categories are defined in \cite{GPS20}, which share a lot of good functorial properties explained in \cite{GPS24S}.\par
Given an exact symplectomorphism $\phi\colon X\to X$ of the Liouville sector $X$, we can define the \emph{mapping cylinder} \[
    M (X,\phi )\coloneqq X\times_{\phi}\mathbb{C}^{\ast}\coloneqq\frac{T^{\ast}\mathbb{R}\times X}{(z,x)\sim (z+1,\phi (x))}
\]which can be regarded as the quotient of the product space $X\times\mathbb{C}$ by a $\mathbb{Z}$-action which acts by $\phi$ on $X$ and by translation by $1$ on $\mathbb{C}$. Its Weinstein structure is induced from $X\times\mathbb{C}$ following definition \ref{defn:G-invariant-Weinstein-sector}, and more explicitly by performing a modification to the product Weinstein structure, using an idea of Kartal \cite{Kartal21}. This modification is done in Proposition \ref{prp:Kartal-modification-Weinstein}. The wrapped Fukaya category of $T^{\ast}\mathbb{R}$ with stops given by preimages of $e^z+e^{-z}$ near infinity was computed in Example \ref{exmp:mirror-to-Coh-G-P1}, and we have
\begin{prp}\label{prp:Fukaya-cat-mapping-cylinder}
    $\mathcal{W} (M(X,\phi),z+1/z)\simeq\left(\mathcal{W} (T^{\ast}\mathbb{R}, e^z+e^{-z})\otimes\mathcal{W} (X)\right)/\mathbb{Z}$.
\end{prp}
\begin{proof}
    By Künneth formula \ref{crl:Kunneth-Fukaya-infinite} combined with the definition of tensor products of $A_{\infty}$-categories in \cite{Ornaghi25}, we otain a quasi-equivalence of $A_{\infty}$-categories \[
        \mathcal{W} (T^{\ast}\mathbb{R},e^z+e^{-z})\otimes\mathcal{W} (X)\hookrightarrow\mathcal{W} (T^{\ast}\mathbb{R}\times X, e^z+e^{-z}).
    \]By theorem \ref{thm:wrapped-Fukaya-quotient}, as we have modified the Weinstein structure on $T^{\ast}\mathbb{R}\times X$ so that it's $\mathbb{Z}$-invariant, we know that \[
        \mathcal{W} (M(X,\phi),z+z^{-1})=\mathcal{W}\left((T^{\ast}\mathbb{R}\times X,e^z+e^{-z})/\mathbb{Z}\right)\simeq\mathcal{W} (T^{\ast}\mathbb{R}\times X,e^z+e^{-z})/\mathbb{Z}
    \]which concludes the Proposition.
% ---- 原文行 380 ----
\end{proof}
On the mirror side, we start with an algebraic variety $Y$ together with a choice of characteristic classes $\xi\in CH^2 (Y)$ that defines a holomorphic line bundle $L_{\xi, Y}$ on $Y$. The subbundle $L_{\xi,Y}^{\times}$ removing the zero section $Y\hookrightarrow L_{\xi, Y}$ is then a principal $\mathbb{G}_m$-bundle over $Y$, and hence carries a free algebraic $\mathbb{G}_m$-action on $Y$. Then we have
\begin{thm}\label{thm:twisted-tensor-P1-bundle}
    Suppose that $Y$ is smooth and quasi-projective, then there is a quasi-equivalence of dg categories\footnote{Here the definition of $\Perf$ is different from the usual definition. See Definition \ref{dfn:dg-enhancement-Coh}.}
$$
\Perf_{\mathbb{G}_m}(\mathbb{P}^1)_{\dg}
\underset{((1),\mathcal{L}(\xi)\otimes-)}{\otimes}
\Perf(Y)_{\dg}
\rightarrow
\Perf(P(Y,\xi))_{\dg}.
$$
\end{thm}

\begin{proof}
    Note that \[
        Y=\left[\frac{L(Y,\xi)^{\times}}{\mathbb{G}_m}\right],
    \]and hence we have quasi-equivalences $\Perf (Y)_{dg}\simeq\Perf_{\mathbb{G}_m} (L(Y,\xi)^{\times})_{dg}$. Note that the action $((1),\mathcal{L} (\xi)\otimes -)$ on $\Perf_{\mathbb{G}_m} (\mathbb{P}^1)_{dg}\otimes\Perf_{\mathbb{G}_m} (L(Y,\xi)^{\times})_{dg}$ is diagonal, by twisted K\"unneth formula we have quasi-embedding \[
        \Perf_{\mathbb{G}_m} (\mathbb{P}^1)\otimes_{\mathbb{Z}=\widehat{\mathbb{G}}_m}\Perf_{\mathbb{G}_m} (L(Y,\xi)^{\times})_{dg}\hookrightarrow\Perf_{\mathbb{G}_m} (\mathbb{P}^1\times L(Y,\xi)^{\times} )_{dg}\cong\Perf (P(Y,\xi))_{dg},
    \]and it suffices to show that this quasi-embedding is essentially surjective. The work of Orlov \cite{Orlov92} showed that in this situation there is a semi-orthogonal decomposition of triangulated categories \[
    D\Perf (P(Y,\xi))_{dg}\simeq\langle D\Perf(Y)_{-1}, D\Perf(Y)_0\rangle,
% ---- 原文行 400 ----
    \]which implies that objects of the form $\pi^{\ast}\mathcal{F}$ and $\pi^{\ast}\mathcal{F}\otimes\mathcal{O}_{P(Y,\xi)} (1)$, where $\mathcal{O}_{P(Y,\xi)} (1)$ is the invertible sheaf whose restriction to each fibre is $\mathcal{O}_{\mathbb{P}^1} (1)$ and $\mathcal{F}\in\Perf (Y)$, generate $\Perf(P(Y,\xi))_{dg}$. Note that these objects are images of the objects $\mathcal{O}_{\mathbb{P}^1} (i)\widetilde{\otimes}\mathcal{F}$ in\\$\Perf_{\mathbb{G}_m} (\mathbb{P}^1)_{dg}\widetilde{\otimes}\Perf (Y)_{dg}$, and therefore the image of the given quasi-embedding generates the target dg category.
\end{proof} 
It's clear from \cite{FLTZ12} and Example \ref{exmp:mirror-to-Coh-G-P1} that $\Perf_{\mathbb{G}_m} (\mathbb{P}^1)_{dg}\simeq\Perf\mathcal{W} (T^{\ast}\mathbb{R},e^z+e^{-z})$, and therefore we have collected enough materials to conclude the proof.
\begin{proof}[Proof of Theorem \ref{thm:HMS-projective-bundle}]
    Since by assumption $\Perf\mathcal{W} (X)^{op}\simeq\Coh (Y)_{dg}=\Perf (Y)_{dg}$ and from the previous discussions $\Perf\mathcal{W} (T^{\ast}\mathbb{R},e^z+e^{-z})^{op}\simeq\Perf_{\mathbb{G}_m} (\mathbb{P}^1)_{dg}$, we have quasi-equivalences between their tensor products \[
        \Perf\mathcal{W} (T^{\ast}\mathbb{R},e^z+e^{-z})^{op}\otimes\Perf\mathcal{W} (X)^{op}\simeq\Perf_{\mathbb{G}_m} (\mathbb{P}^1)_{dg}\otimes\Perf (Y)_{dg}.
    \]On the $A$-model the $\mathbb{Z}$-action on the product is provided by shifting by $1$ on $T^{\ast}\mathbb{R}$ and by the symplectomorphism $\phi$ on $\mathcal{W} (X)$, which coincides with the categorical $\mathbb{Z}$-action on the $B$-side coming from the $\mathbb{G}_m$-action. Therefore after taking $\mathbb{Z}$-orbit categories and combining with Proposition \ref{prp:Fukaya-cat-mapping-cylinder} and Theorem \ref{thm:twisted-tensor-P1-bundle} we conclude the proof.
\end{proof}
\begin{proof}[Proof of Corollary \ref{crl:mirror-symmetry-line-principal-bundles}]
    It suffices to show that the inclusion functor $i_0\colon\Coh (Y)\to\Coh (P(Y,\xi))$ is mirror to the Orlov functor $\cup_0\colon\mathcal{W} (X)\to\mathcal{W} (M(X,\phi),z+z^{-1})$. Note that by construction, $i_0$ can be factored as \[
        \Perf (Y)_{dg}\xrightarrow{i_0}\Perf (Y)_{dg}\otimes\Perf_{\mathbb{G}_m}(\mathbb{P}^1)_{dg}\to\Perf (P(Y,\xi))_{dg},
    \]and similarly $\cup_0$ can be factored as \[
        \Perf\mathcal{W} (X)^{op}\xrightarrow{\cup_0}\Perf\mathcal{W} (X\times (T^{\ast}\mathbb{R},e^z+e^{-z}))^{op}\to\Perf\mathcal{W} (M(X,\phi),z+z^{-1})^{op},
    \]where the second map are all quotient by the $\mathbb{Z}$-action, which we know to be mirror to each other. On the $B$-side the functor $i_0$ is given by sending any $\mathcal{F}\in\Coh (Y)$ to the sheaf $\mathcal{F}\otimes\mathcal{O}_0 [0]$, while on the $A$-side $\cup_0$ sends $L\in\mathcal{W} (X)$ to $L\otimes D[0]$, where $D[0]\in\Perf\mathcal{W} (T^{\ast}\mathbb{R},e^z+e^{-z})^{op}\simeq\Perf_{\mathbb{G}_m} (\mathbb{P}^1)_{dg}$ is the linking disk to the stop labelled as $0$, which up to a shift is identified with $\mathcal{O}_0 [0]$. Therefore $i_0$ and $\cup_0$ are mirror to each other and by taking localizations we conclude the proof. The proof for principal bundles is basically the same, with $0$ replaced by $\infty$.
\end{proof}
\subsection{Sectorial and proper descents.} Now we move forward to the proof of the theorems \ref{crl:gluing-mirror-symmetry} and \ref{crl:gluing-mirror-symmetry-circ}. Although stated in a pretty general way, it suffices for us to understand the case when $k=2$, which will follow from the discussion of Example \ref{exmp:wrapped-Fukaya-cat-twice-punctured-equivariant}.\par
Given $B$-model data $(Y,\xi_1,\xi_2)$, we construct the space \[
    \mathcal{Y} (\xi_1,\xi_2)=P(Y,\xi_1)\cup_Y P(Y,\xi_2)
\]by gluing the two projective bundles by identifying the section at $0$ and the section at $\infty$. Note that there is a $\mathbb{G}_m^2$-action on $\mathcal{Y} (\xi_1,\xi_2)$, given by the $\mathbb{G}_m$-actions on each $P(Y,\xi_i)$ that extend trivially to the other component. From the proper pushforward diagram \[
    \begin{tikzcd}
% ---- 原文行 420 ----
        Y\arrow[r,"i_0"]\arrow[d,"i_{\infty}"]&P(Y,\xi_2)\arrow[d]\\ 
        P(Y,\xi_1)\arrow[r]&\mathcal{Y} (\xi_1,\xi_2)
    \end{tikzcd}
\]we get directly the proper pushforward diagram of quotient stacks\footnote{This follows directly from the definition of quotient stacks and valuative criterion. See e.g. \cite{Alper26}.} \[
    \begin{tikzcd}
        {[Y/\mathbb{G}_m^2]}\arrow[r,"i_0"]\arrow[d,"i_{\infty}"]&{\big[P(Y,\xi_2)/\mathbb{G}_m^2\big]}\arrow[d]\\ 
        \big[P(Y,\xi_1)/\mathbb{G}_m^2\big]\arrow[r]&\big[\mathcal{Y}(\xi_1,\xi_2)/\mathbb{G}_m^2\big]
    \end{tikzcd}
\]The homotopy pushforward formula for quasi-coherent sheaves for pushforward diagrams of closed immersions, which follows from Gaitsgory-Rozenblyum \cite[Theorem A.1.2]{GR17II} together with descent property for quasi-coherent sheaves, implies the homotopy pushforward diagram \[
    \begin{tikzcd}
        {\IndCoh [Y/\mathbb{G}_m^2]_{\infty}}\arrow[r,"i_0"]\arrow[d,"i_{\infty}"]&{\IndCoh\big[P(Y,\xi_2)/\mathbb{G}_m^2\big]_{\infty}}\arrow[d]\\ 
        \IndCoh\big[P(Y,\xi_1)/\mathbb{G}_m^2\big]_{\infty}\arrow[r]&\IndCoh\big[\mathcal{Y}(\xi_1,\xi_2)/\mathbb{G}_m^2\big]_{\infty}
    \end{tikzcd}
\]and since coherent sheaves are pushed to coherent sheaves via closed immersions, it follows that the same diagram for coherent sheaves is also a homotopy pushforward diagram. Using the quasi-equivalence between $dg$- and $\infty$-categories proved in \cite{Cohn16}, we have
\begin{prp}\label{prp:pushforward-coherent-stack}
    There is a quasi-equivalence of $A_{\infty}$-categories \[
        \Coh_{\mathbb{G}_m^2}\mathcal{Y} (\xi_1,\xi_2)_{dg}\simeq\hocolim\left(\Perf_{\mathbb{G}_m^2} (P(Y,\xi_1))_{dg}\leftarrow\Perf_{\mathbb{G}_m^2} (Y)_{dg}\rightarrow\Perf_{\mathbb{G}_m^2} (P(Y,\xi_2))_{dg}\right).
    \]
\end{prp}
Now we turn to the mirror. Since the extra $\mathbb{G}_m$-action is also identity on $Y$, we have that $[P(Y,\xi_i)/\mathbb{G}_m^2]=[\mathbb{P}^1/\mathbb{G}_m]\times Y\times [\mathrm{pt}/\mathbb{G}_m]$. As $\Perf_{\mathbb{G}_m} (\mathbb{P}^1)_{dg}\simeq\Perf\mathcal{W} (T^{\ast}\mathbb{R},e^z+e^{-z})$, it suffices to understand Homological Mirror Symmetry for $[\mathrm{pt}/\mathbb{G}_m]$.
% ---- 原文行 440 ----
\begin{lmm}
    There is a quasi-equivalence of $A_{\infty}$-categories \[
        \Perf\mathcal{W}\left(\bigsqcup_{i\in\mathbb{Z}}T^{\ast}[0,1]\right)^{op}\simeq\Perf_{\mathbb{G}_m} (\mathrm{pt})_{dg}.
    \]
\end{lmm}
\begin{proof}
    Note that the left-hand side is the homotopy coproduct \[
       \Perf \mathcal{W}\left(\bigsqcup_{i\in\mathbb{Z}}T^{\ast} [0,1]\right)^{op}\simeq\coprod_{i\in\mathbb{Z}}\Perf\mathcal{W} (T^{\ast} [0,1])^{op}\simeq \coprod_{i\in\mathbb{Z}}\Perf (\mathrm{pt})_{dg}\simeq\coprod_{i\in\mathbb{Z}}\Perf\field\text{-}\Mod_{dg},
    \]while on the other hand, from \cite[Proposition 4.9]{TV18} we know that the $\mathbb{G}_m$-equivariant quasi-coherent category $\QCoh_{\mathbb{G}_m} (\mathrm{pt})_{dg}\simeq\Rep_{\Bbbk} (\mathbb{G}_m)_{dg}$ is equivalent to the category of representations of $\mathbb{G}_m$, and since a linear $\mathbb{G}_m$-representation is the same as a $\mathbb{Z}$-grading, we see that $\Rep_{\Bbbk} (\mathbb{G}_m)_{dg}$ is equivalent to category of complexes of $\field$-vector spaces equipped with an extra $\mathbb{Z}$-grading. By taking compact objects, we know that \[
        \Perf_{\mathbb{G}_m} (\mathrm{pt})_{\infty}\simeq\left(\Rep_{\Bbbk} (\mathbb{G}_m)_{\infty}\right)^{\omega}\simeq\coprod_{i\in\mathbb{Z}}\Perf\Bbbk\text{-}\Mod_{\infty}
    \]using the fact that $\Perf_{\mathbb{G}_m} (\mathrm{pt})_{dg}$ is the compact subcategory of $\QCoh_{\mathbb{G}_m} (\mathrm{pt})_{dg}$ \cite{BZFN10}. We then use the equivalence between $dg$-categories and $\Bbbk$-linear stable $\infty$-categories \cite{Cohn16} to conclude the proof.
\end{proof}
From this lemma and the previous discussions, we get a quasi-equivalence of $A_{\infty}$-categories \begin{align}\label{eq:equivariant-mirror-Gm2}
    \begin{split}
    \Perf\mathcal{W}\left(\bigsqcup_{i\in\mathbb{Z}} (T^{\ast}\mathbb{R},e^z+e^{-z})\right)^{op}&\simeq\Perf\mathcal{W} (T^{\ast}\mathbb{R},e^z+e^{-z})^{op}\otimes\Perf_{\mathbb{G}_m} (\mathrm{pt})\\ 
    &\simeq\Perf_{\mathbb{G}_m} (\mathbb{P}^1)\otimes\Perf_{\mathbb{G}_m} (\mathrm{pt})\simeq\Perf_{\mathbb{G}_m^2} (\mathbb{P}^1).
    \end{split}
\end{align}
\begin{proof}[Proof of Corollary \ref{crl:gluing-mirror-symmetry}]
    As already discussed, it suffices for us to verify \[
% ---- 原文行 460 ----
        \Perf\mathcal{W}\left(\mathcal{X} (\phi_1,\phi_2)\right)^{op}\simeq\Coh\left(\mathcal{Y}(\xi_1,\xi_2)\right)
    \]for chosen mirror data $(X,\phi_1,\phi_2)$ and $(Y,\xi_1,\xi_2)$. Equivalence \eqref{eq:equivariant-mirror-Gm2} together with analyse of mirror functors provides a quasi-equivalence between the diagrams \[
\begin{aligned}
&\Perf\left(
\begin{tikzcd}
\mathcal{W}\left(\bigsqcup_{\mathbb{Z}^2}T^{\ast}[0,1]\right)
  \arrow[r]\arrow[d]
&
\mathcal{W}\left(\bigsqcup_{\mathbb{Z}}(\mathbb{C},e^{z_1}+e^{-z_1})\right)
\\
\mathcal{W}\left(\bigsqcup_{\mathbb{Z}}(\mathbb{C},e^{z_2}+e^{-z_2})\right)
&
\end{tikzcd}
\right)^{op}
\\[1em]
&\simeq
\left(
\begin{tikzcd}
\Coh[\mathrm{pt}/\mathbb{G}_m^2]_{dg}\arrow[r]\arrow[d]
&
% ---- 原文行 480 ----
\Coh[\mathbb{P}^1/\mathbb{G}_m^2]_{dg}
\\
\Coh[\mathbb{P}^1/\mathbb{G}_m^2]_{dg}
&
\end{tikzcd}
\right).
\end{aligned}
\]and therefore we have $\Perf\mathcal{W} (\mathcal{X})^{op}\simeq\Coh_{\mathbb{G}_m^2} (\mathbb{P}^1\cup_{0\sim\infty}\mathbb{P}^1)$, which provides the desired categorical equivalence by tensoring with $\Perf\mathcal{W} (X)^{op}\simeq\Coh (Y)_{dg}$ and then quotient by the canonical $\mathbb{Z}^2$-action.
\end{proof}
\begin{proof}[Proof of Corollary \ref{crl:gluing-mirror-symmetry-circ}]
    The proof is completely parallel to the proof of Corollary \ref{crl:gluing-mirror-symmetry}.
\end{proof}
\subsection{Toric bundles.} Note from equivariant coherent-constructible correspondence \cite{FLTZ11} that we have a quasi-equivalence of dg categories \[
    \Sh_{\widetilde{\Lambda}_{\Sigma}} (\mathbb{R}^n)^{cc}\simeq\Coh_{\mathbb{G}_m^n} (X_{\Sigma})
\]when $X_{\Sigma}$ is a smooth toric variety.
Here $\Sh_{\widetilde{\Lambda}_{\Sigma}} (\mathbb{R}^n)^{cc}$ is the dg category of constructible complexes of sheaves with compact supports, and by \cite[Lemma 4.18]{GPS24M} we have \[
    \Sh_{\widetilde{\Lambda}_{\Sigma}} (\mathbb{R}^n)^{cc}\simeq\Sh_{\widetilde{\Lambda}_{\Sigma}} (\mathbb{R}^n)^{\omega},
\]which we know from sheaf-Floer correspondence \ref{prp:inf-type-sheaf-Floer-Rn} to be quasi-equivalent to the wrapped Fukaya category $\Perf\mathcal{W} (T^{\ast}\mathbb{R}^n,\widetilde{\Lambda}_{\Sigma})^{op}$.
Therefore we have a quasi-equivalence of $A_{\infty}$-categories \[
    \Perf\mathcal{W} (T^{\ast}\mathbb{R}^n,\widetilde{\Lambda}_{\Sigma})^{op}\simeq\Coh_{\mathbb{G}_m^n} (X_{\Sigma})
% ---- 原文行 500 ----
\]whenever $X_{\Sigma}$ is smooth, or equivalently $\Sigma$ is an integral convex polytope.
\begin{proof}[Proof of Corollary \ref{crl:toric-bundle-HMS}]
    From $A$-side we have \[
        \mathcal{W}\left[\mathcal{X}_{\Sigma} (X;\phi_1,\dotsb ,\phi_k)\right]\simeq\mathcal{W} (T^{\ast}\mathbb{R}^k,\widetilde{\Lambda}_{\Sigma})\otimes_{\mathbb{Z}^k}\mathcal{W} (X),
    \]while on the $B$-side we can write $\mathcal{Y}_{\Sigma} (Y;\xi_1,\dotsb ,\xi_k)=\left(X_{\Sigma}\times \prod_{i=1}^k L(Y,\xi_i)^{\times}\right)/\mathbb{G}_m^k$ and so \[
        \Coh\left[\mathcal{Y} (Y;\xi_1,\dotsb ,\xi_k)\right]\simeq\left[\Coh_{\mathbb{G}_m^k} (X_{\Sigma})\otimes\Coh_{\mathbb{G}_m^k} \left(\prod_{i=1}^kL(Y,\xi_i)^{\times}\right)\right]\Big/\mathbb{G}_m^k.
    \]Pairing these two equivalences concludes.
\end{proof}
\subsection{Hirzebruch surfaces.} The last part of the proof is contributed to Homological Mirror Symmetry for Hirzebruch surfaces. Note that the HMS part of Corollary \ref{crl:mirror-to-Hirzebruch} is a direct Corollary of Theorem \ref{thm:HMS-projective-bundle}, and the only non-trivial part is the geometric identification of the constructed space $\rF_n$.
\begin{proof}[Proof of Corollary \ref{crl:mirror-to-Hirzebruch}]
     Consider the affine hypersurface \[
            Z\coloneqq\{(x,y,z)\in\mathbb{C}^3\vert x^{n+1}+yz=0\}\subseteq\mathbb{C}^3
        \]which admits a holomorphic map \[
            \pi\colon Z\to\mathbb{C},\ (x,y,z)\mapsto x
        \]so that $\pi^{-1} (0)=\{(y,z)\vert yz=0\}$ is singular nodal while $\pi^{-1} (t)=\{(y,z)\vert yz =t^{n+1}\}$ is smooth affine, which from its expression (and Riemann-Hurwitz formula) is clearly isomorphic to $\mathbb{C}^{\ast}$. We can also verify that the monodromy action around the singular fibre $x=0$ is the same as the Dehn twist $\phi^n$, and therefore we have $\rF_n\cong_{alg} Z\setminus\{x=0\}$. Note that $\rF_n$ admits a projection $\rF_n\to\mathbb{C}_x^{\ast}\times\mathbb{C}_y^{\ast}$ where the fibre at each pair $(x,y)$ consists exactly of a single point, and therefore we have $\mathcal{X}\cong_{alg}\mathbb{C}_x^{\ast}\times\mathbb{C}_y^{\ast}$. This isomorphism does not change the coordinates of $x$ and $y$, so the superpotential under this algebraic isomorphism is still $x+x^{-1} +y+y^{-1}$.
\end{proof}
